\documentclass[11pt]{article}

\usepackage[utf8]{inputenc}
\usepackage[T1]{fontenc}
\usepackage{geometry}
\usepackage{amsmath}
\usepackage{amssymb}
\usepackage{booktabs}
\usepackage{hyperref}
\usepackage{xurl}
\usepackage{xcolor}
\usepackage{listings}
\usepackage{enumitem}
\usepackage{parskip}
\usepackage{graphicx}
\graphicspath{{../images/}}

\geometry{margin=1in}

\lstset{
  language=Java,
  basicstyle=\ttfamily\small,
  keywordstyle=\color{blue},
  commentstyle=\color{gray},
  stringstyle=\color{teal},
  showstringspaces=false,
  columns=fullflexible,
  frame=single,
  framerule=0.2pt,
  breaklines=true
}

\setlist[itemize]{topsep=0.3em,itemsep=0.2em}
\setlist[enumerate]{topsep=0.3em,itemsep=0.2em}

% Simple per-section source line.
\newcommand{\notesources}[1]{\par\noindent{\small\color{gray}\textit{Sources:} \sloppy #1\par}}

% Unit-wise source strings for this subject (used by slides/notes).
% Path is relative to the lecture `latex/` folder (the compilation CWD).
% OOPS with Java (CSEG1044) — unit-wise sources
%
% These are short, student-facing reference strings intended to be used with:
%   - \slidesources{...}  (in beamer slides)
%   - \notesources{...}   (in lecture notes)
%
% Style inspiration: other subjects in My-Subjects/ use semicolon-separated sources
% and include an "accessed" date for web documentation.

\newcommand{\OOPSUnitISources}{Schildt, \textit{Java: A Beginner's Guide} (10e), McGraw-Hill (Java fundamentals + OOP basics).; Oracle Java Tutorials: Getting Started + Language Basics + Classes and Objects (accessed \today).; Java SE API docs (e.g., \texttt{String}, \texttt{Scanner}) (accessed \today).}

\newcommand{\OOPSUnitIISources}{Schildt, \textit{Java: The Complete Reference} (12e), McGraw-Hill (classes, inheritance, packages, interfaces).; Bloch, \textit{Effective Java} (3e), Addison-Wesley, 2018 (design choices; interfaces vs abstract classes).; Oracle Java Tutorials: Classes and Objects + Interfaces + Packages (accessed \today).; Java SE API docs (e.g., \texttt{Comparable}, \texttt{Runnable}) (accessed \today).}

\newcommand{\OOPSUnitIIISources}{Schildt, \textit{Java: The Complete Reference} (12e) (nested/inner classes, exceptions, threads, I/O).; Oracle Java Tutorials: Nested Classes + Exceptions + Concurrency + Basic I/O (accessed \today).; Java SE API docs (e.g., \texttt{Thread}, \texttt{AutoCloseable}) (accessed \today).}

\newcommand{\OOPSUnitIVSources}{Schildt, \textit{Java: The Complete Reference} (12e) (generics, lambdas, Swing, JDBC).; Bloch, \textit{Effective Java} (3e) (generics + lambdas).; Oracle Java Tutorials: Generics + Lambda Expressions + Swing + JDBC Basics (accessed \today).; Java SE API docs (e.g., \texttt{java.util.function}, \texttt{javax.swing}, \texttt{java.sql}) (accessed \today).}

\newcommand{\OOPSUnitVSources}{Schildt, \textit{Java: The Complete Reference} (12e) (Collections Framework + wrapper classes).; Oracle Java docs: Collections Framework overview (accessed \today).; Java SE API docs (\texttt{java.util} collections; wrapper classes in \texttt{java.lang}) (accessed \today).}

\newcommand{\OOPSUnitVISources}{Git SCM: \textit{Pro Git} (2e) (version control workflow), accessed \today.; GitHub Docs (repositories, issues, pull requests), accessed \today.; Oracle Java Tutorials: Swing + JDBC (accessed \today).; JUnit 5 User Guide (accessed \today).; Apache Maven Project: Getting Started (accessed \today).; Gradle User Manual: Getting Started (accessed \today).; UML class diagrams: UML 2.x overview/spec (accessed \today).}

\newcommand{\OOPSMiscWhyInterfacesSources}{Schildt, \textit{Java: The Complete Reference} (12e) (interfaces).; Bloch, \textit{Effective Java} (3e) (prefer interfaces where appropriate).; Oracle Java Tutorials: Interfaces (accessed \today).; Java SE API docs (e.g., \texttt{Comparable}, \texttt{Runnable}) (accessed \today).}



\title{Object Oriented Programming (CSEG1044)\\Unit 04 -- Lecture 05 Notes\\Swing Event Handling + JDBC (Persistence)}
\author{Tofik Ali}
\date{\today}

\begin{document}
\maketitle
\tableofcontents

\section{Lecture Overview}
In the last lecture we built a Swing UI skeleton.
Now we make it interactive and persistent:
\begin{itemize}
  \item \textbf{Events}: a button click triggers code (save/clear/exit).
  \item \textbf{Persistence}: data is stored in a database using JDBC so it is not lost after closing the app.
\end{itemize}
\notesources{\OOPSUnitIVSources}

\section{Syllabus Alignment (Unit 04)}
\textbf{Unit 04:} Swing event handling; database connectivity using JDBC; basic CRUD

\section{Learning Outcomes}
After this lecture, you should be able to:
\begin{itemize}
  \item Register Swing listeners (e.g., \texttt{ActionListener}) and handle common events safely.
  \item Explain JDBC flow: connect $\rightarrow$ prepare statement $\rightarrow$ execute $\rightarrow$ read \texttt{ResultSet}.
  \item Differentiate \texttt{Statement}, \texttt{PreparedStatement}, \texttt{CallableStatement}.
  \item Implement basic CRUD operations and map rows to objects.
\end{itemize}

\section{Key Terms (Quick)}
\begin{itemize}
  \item \textbf{Event source}: UI component that triggers an event (button, menu item).
  \item \textbf{Event object}: information about the event (e.g., \texttt{ActionEvent}).
  \item \textbf{Listener}: object/function that handles events (e.g., \texttt{ActionListener}).
  \item \textbf{EDT}: Event Dispatch Thread (Swing UI thread).
  \item \textbf{JDBC}: Java Database Connectivity API.
  \item \textbf{Connection}: active DB connection.
  \item \textbf{PreparedStatement}: parameterized SQL statement (safer and common).
  \item \textbf{ResultSet}: table-like cursor over query results.
  \item \textbf{CRUD}: Create, Read, Update, Delete.
  \item \textbf{DAO}: Data Access Object (separate DB code from UI).
\end{itemize}

\section{Detailed Notes}
\subsection{Swing event model (listeners)}
Swing uses a listener-based model:
\begin{itemize}
  \item You create a component (e.g., \texttt{JButton saveBtn}).
  \item You register a listener: \texttt{saveBtn.addActionListener(...)}.
  \item When user clicks, Swing calls your listener method on the EDT.
\end{itemize}

\textbf{Common controls and events:}
\begin{itemize}
  \item \texttt{JButton} $\rightarrow$ \texttt{ActionListener}
  \item \texttt{JTextField} $\rightarrow$ \texttt{ActionListener} (press Enter), \texttt{DocumentListener} (text changes)
  \item \texttt{JTable} $\rightarrow$ selection listeners (advanced)
\end{itemize}

\paragraph{Important rule.} Do not run slow DB operations on the EDT; it can freeze the UI.
In a basic course project, keep operations small or use a background thread/SwingWorker (advanced).

\subsection{JDBC core flow (pattern you should memorize)}
\begin{enumerate}
  \item Load/ensure driver is available (often automatic in modern JDBC).
  \item Get a connection: \texttt{DriverManager.getConnection(url, user, pass)}.
  \item Create a statement (usually \texttt{PreparedStatement}).
  \item Execute:
  \begin{itemize}
    \item \texttt{executeUpdate()} for INSERT/UPDATE/DELETE
    \item \texttt{executeQuery()} for SELECT
  \end{itemize}
  \item Read results using \texttt{ResultSet} (cursor).
  \item Close resources (use \texttt{try-with-resources}).
\end{enumerate}

\subsection{\texttt{Statement} vs \texttt{PreparedStatement} vs \texttt{CallableStatement}}
\begin{itemize}
  \item \texttt{Statement}: simplest; builds SQL as a string (risk of SQL injection; avoid for user input).
  \item \texttt{PreparedStatement}: parameterized with \texttt{?}; safer and faster for repeated queries.
  \item \texttt{CallableStatement}: calls stored procedures (depends on DB).
\end{itemize}

\subsection{\texttt{ResultSet} traversal}
\begin{itemize}
  \item \texttt{while (rs.next()) \{ ... \}} moves row by row.
  \item Read columns by label or index: \texttt{rs.getInt("id")}, \texttt{rs.getString("name")}.
  \item Map each row to an object (e.g., \texttt{Resource}).
\end{itemize}

\subsection{Persistent data design (UI vs DAO)}
Avoid writing SQL inside button click code.
Preferred structure:
\begin{itemize}
  \item UI layer: reads fields, validates, shows messages.
  \item DAO layer: contains SQL and JDBC code.
  \item Model layer: \texttt{Resource} class that represents one row/entity.
\end{itemize}
\notesources{\OOPSUnitIVSources}

\section{Worked Example: Save a resource from Swing to DB (skeleton)}
\subsection*{1) Model class}
\begin{lstlisting}
public class Resource {
    private final String name;
    private final String type;
    private final int capacity;

    public Resource(String name, String type, int capacity) {
        this.name = name;
        this.type = type;
        this.capacity = capacity;
    }

    public String getName() { return name; }
    public String getType() { return type; }
    public int getCapacity() { return capacity; }
}
\end{lstlisting}

\subsection*{2) DAO (insert + select)}
\begin{lstlisting}
import java.sql.*;
import java.util.ArrayList;
import java.util.List;

public class ResourceDao {
    private final String url = "jdbc:mysql://localhost:3306/crm"; // example
    private final String user = "root";
    private final String pass = "password";

    public void insert(Resource r) throws SQLException {
        String sql = "INSERT INTO resources(name, type, capacity) VALUES(?,?,?)";
        try (Connection con = DriverManager.getConnection(url, user, pass);
             PreparedStatement ps = con.prepareStatement(sql)) {
            ps.setString(1, r.getName());
            ps.setString(2, r.getType());
            ps.setInt(3, r.getCapacity());
            ps.executeUpdate();
        }
    }

    public List<Resource> findAll() throws SQLException {
        String sql = "SELECT name, type, capacity FROM resources";
        List<Resource> out = new ArrayList<>();
        try (Connection con = DriverManager.getConnection(url, user, pass);
             PreparedStatement ps = con.prepareStatement(sql);
             ResultSet rs = ps.executeQuery()) {
            while (rs.next()) {
                out.add(new Resource(
                    rs.getString("name"),
                    rs.getString("type"),
                    rs.getInt("capacity")
                ));
            }
        }
        return out;
    }
}
\end{lstlisting}

\subsection*{3) Event handler (button click)}
\begin{lstlisting}
saveBtn.addActionListener(e -> {
    String name = nameField.getText().trim();
    String type = typeField.getText().trim();
    String capText = capacityField.getText().trim();

    if (name.isEmpty() || type.isEmpty() || capText.isEmpty()) {
        JOptionPane.showMessageDialog(f, "All fields required");
        return;
    }

    int cap;
    try {
        cap = Integer.parseInt(capText);
        if (cap <= 0) throw new NumberFormatException();
    } catch (NumberFormatException ex) {
        JOptionPane.showMessageDialog(f, "Capacity must be a positive integer");
        return;
    }

    try {
        dao.insert(new Resource(name, type, cap));
        JOptionPane.showMessageDialog(f, "Saved!");
    } catch (SQLException ex) {
        JOptionPane.showMessageDialog(f, "DB error: " + ex.getMessage());
    }
});
\end{lstlisting}

\section{In-Class Exercises (with brief solutions)}
\begin{enumerate}
  \item Which is safer for user input: \texttt{Statement} or \texttt{PreparedStatement}?\par
  \textbf{Answer:} \texttt{PreparedStatement}.

  \item Why should heavy DB operations not run on the EDT?\par
  \textbf{Answer:} the UI can freeze because Swing processes events on the EDT.

  \item What method is used for SELECT queries: \texttt{executeQuery()} or \texttt{executeUpdate()}?\par
  \textbf{Answer:} \texttt{executeQuery()}.
\end{enumerate}

\section{Self-Study Practice (no solutions)}
\begin{enumerate}
  \item Add an \texttt{updateCapacity(name, newCapacity)} method in \texttt{ResourceDao}.
  \item Add a ``Load All'' button that displays all resources in a text area/table.
  \item Replace DB credentials with a config file and load them at startup (optional).
\end{enumerate}

\section{Lab Alignment (recommended)}
\begin{itemize}
  \item Exp-16: GUI using Swing
  \item Exp-17: Database connectivity using JDBC
\end{itemize}

\section{Checkpoint Questions (with sample answers)}
\textbf{Q1.} Why is \texttt{PreparedStatement} preferred?\par
\textbf{Sample answer:} It allows parameterized queries, prevents SQL injection for user input, and can be more efficient for repeated execution.

\vspace{0.6em}
\textbf{Q2.} What is the role of a DAO?\par
\textbf{Sample answer:} A DAO separates database code from UI/business logic, making the application easier to maintain, test, and change.

\end{document}
